% math-equations.tex — a tour of LaTeX mathematics: inline and display,
% numbered and aligned, matrices, cases, integrals, and theorem environments.
% Compile with: pdflatex math-equations.tex
%
% The results stated here are standard textbook material restated in the
% author's own words; nothing is presented as new.

\documentclass[11pt,a4paper]{article}

\usepackage[utf8]{inputenc}
\usepackage[T1]{fontenc}
\usepackage{lmodern}
\usepackage{amsmath}
\usepackage{amssymb}
\usepackage{amsthm}
\usepackage{geometry}

\geometry{margin=2.6cm}

\theoremstyle{plain}
\newtheorem{theorem}{Theorem}[section]
\newtheorem{lemma}[theorem]{Lemma}
\newtheorem{corollary}[theorem]{Corollary}

\theoremstyle{definition}
\newtheorem{definition}[theorem]{Definition}
\newtheorem{example}[theorem]{Example}

\theoremstyle{remark}
\newtheorem*{remark}{Remark}

\newcommand{\R}{\mathbb{R}}
\newcommand{\N}{\mathbb{N}}
\newcommand{\C}{\mathbb{C}}
\newcommand{\norm}[1]{\left\lVert #1 \right\rVert}
\newcommand{\abs}[1]{\left\lvert #1 \right\rvert}
\DeclareMathOperator{\tr}{tr}
\DeclareMathOperator*{\argmin}{arg\,min}

\title{A Working Set of Mathematical Notation}
\author{Sample Assets}
\date{}

\begin{document}

\maketitle

\section{Inline and display}

An inline formula sits in the line of text, like $e^{i\pi} + 1 = 0$, and should
not disturb the leading. A displayed formula is set apart:
\[
  \sum_{n=1}^{\infty} \frac{1}{n^{2}} = \frac{\pi^{2}}{6}.
\]
A numbered display can be referred to later:
\begin{equation}
  \label{eq:gaussian}
  \int_{-\infty}^{\infty} e^{-x^{2}} \, dx = \sqrt{\pi}.
\end{equation}
Equation~\eqref{eq:gaussian} is the one identity every typesetting sample
eventually contains.

\section{Alignment}

Several steps of a derivation align on the relation symbol:
\begin{align}
  (a + b)^{3} &= (a + b)(a + b)^{2} \notag \\
              &= (a + b)(a^{2} + 2ab + b^{2}) \notag \\
              &= a^{3} + 3a^{2}b + 3ab^{2} + b^{3}.
  \label{eq:cube}
\end{align}
Two unrelated equations can share a display without sharing an alignment point:
\begin{gather}
  \nabla \cdot \mathbf{E} = \frac{\rho}{\varepsilon_{0}}, \\
  \nabla \times \mathbf{B} - \mu_{0}\varepsilon_{0}
    \frac{\partial \mathbf{E}}{\partial t} = \mu_{0}\mathbf{J}.
\end{gather}

\section{Matrices and arrays}

\[
  A =
  \begin{pmatrix}
    a_{11} & a_{12} & \cdots & a_{1n} \\
    a_{21} & a_{22} & \cdots & a_{2n} \\
    \vdots & \vdots & \ddots & \vdots \\
    a_{m1} & a_{m2} & \cdots & a_{mn}
  \end{pmatrix},
  \qquad
  \det
  \begin{vmatrix}
    a & b \\
    c & d
  \end{vmatrix}
  = ad - bc .
\]

For a square matrix $A \in \R^{n \times n}$ with eigenvalues
$\lambda_{1}, \dots, \lambda_{n}$,
\[
  \tr(A) = \sum_{i=1}^{n} \lambda_{i},
  \qquad
  \det(A) = \prod_{i=1}^{n} \lambda_{i}.
\]

\section{Cases and piecewise definitions}

\begin{definition}[Sign function]
For $x \in \R$,
\[
  \operatorname{sgn}(x) =
  \begin{cases}
    -1 & \text{if } x < 0, \\
     0 & \text{if } x = 0, \\
     1 & \text{if } x > 0.
  \end{cases}
\]
\end{definition}

\section{Limits, sums, and products}

\[
  \lim_{n \to \infty} \left( 1 + \frac{x}{n} \right)^{n} = e^{x},
  \qquad
  \prod_{k=1}^{n} k = n!,
  \qquad
  \binom{n}{k} = \frac{n!}{k!\,(n-k)!}.
\]

The derivative of a product and the chain rule, side by side:
\begin{align*}
  \frac{d}{dx}\bigl( f(x)g(x) \bigr) &= f'(x)g(x) + f(x)g'(x), \\
  \frac{d}{dx} f\bigl( g(x) \bigr)   &= f'\bigl( g(x) \bigr)\, g'(x).
\end{align*}

\section{Theorems}

\begin{theorem}[Triangle inequality]
\label{thm:triangle}
For all $u, v$ in a normed vector space,
\[
  \norm{u + v} \le \norm{u} + \norm{v}.
\]
\end{theorem}

\begin{proof}
Expand the square of the left-hand side and apply the Cauchy--Schwarz
inequality:
\[
  \norm{u + v}^{2}
    = \norm{u}^{2} + 2\langle u, v \rangle + \norm{v}^{2}
    \le \norm{u}^{2} + 2\norm{u}\norm{v} + \norm{v}^{2}
    = \bigl( \norm{u} + \norm{v} \bigr)^{2}.
\]
Both sides are non-negative, so taking square roots preserves the inequality.
\end{proof}

\begin{lemma}
\label{lem:reverse}
Under the hypotheses of Theorem~\ref{thm:triangle},
$\bigl\lvert \norm{u} - \norm{v} \bigr\rvert \le \norm{u - v}$.
\end{lemma}

\begin{corollary}
The norm $\norm{\cdot}$ is continuous.
\end{corollary}

\begin{remark}
Lemma~\ref{lem:reverse} is usually called the reverse triangle inequality, and
follows from Theorem~\ref{thm:triangle} by writing $u = (u - v) + v$.
\end{remark}

\begin{example}
Let $f(x) = \abs{x}$ on $\R$. Then $f$ is continuous everywhere and
differentiable everywhere except at $x = 0$, where the left and right difference
quotients converge to $-1$ and $1$ respectively.
\end{example}

\section{An optimisation problem}

A least-squares fit written as a minimisation:
\begin{equation}
  \hat{\beta} = \argmin_{\beta \in \R^{p}}
    \; \norm{y - X\beta}_{2}^{2} + \lambda \norm{\beta}_{2}^{2},
  \qquad \lambda \ge 0,
\end{equation}
with the closed-form solution
\[
  \hat{\beta} = \left( X^{\top}X + \lambda I \right)^{-1} X^{\top} y ,
\]
which exists for every $\lambda > 0$ even when $X^{\top}X$ is singular.

\section{Symbols in running text}

Set membership $x \in A$, subset $A \subseteq B$, union $A \cup B$, intersection
$A \cap B$, and complement $A^{c}$. Logical quantifiers
$\forall \varepsilon > 0 \; \exists \delta > 0$, implication
$P \Rightarrow Q$, equivalence $P \iff Q$. Greek letters
$\alpha, \beta, \gamma, \delta, \varepsilon, \theta, \lambda, \mu, \sigma,
\varphi, \omega$ and their capitals $\Gamma, \Delta, \Theta, \Lambda, \Sigma,
\Phi, \Omega$. Number systems $\N \subset \mathbb{Z} \subset \mathbb{Q}
\subset \R \subset \C$.

\end{document}
